%%%%%%%%%%%%%%%%%%%%%%%%%%%%%%%%%%%%%%%%%
% Masters/Doctoral Thesis
% LaTeX Template
% Version 1.43 (17/5/14)
%
% This template has been downloaded from:
% http://www.LaTeXTemplates.com
%
% Original authors:
% Steven Gunn
% http://users.ecs.soton.ac.uk/srg/softwaretools/document/templates/
% and
% Sunil Patel
% http://www.sunilpatel.co.uk/thesis-template/
%
% License:
% CC BY-NC-SA 3.0 (http://creativecommons.org/licenses/by-nc-sa/3.0/)
%
% Note:
% Make sure to edit document variables in the Thesis.cls file
%
%%%%%%%%%%%%%%%%%%%%%%%%%%%%%%%%%%%%%%%%%

%----------------------------------------------------------------------------------------
%	PACKAGES AND OTHER DOCUMENT CONFIGURATIONS
%----------------------------------------------------------------------------------------

\documentclass[11pt, oneside]{Thesis} % The default font size and one-sided printing (no margin offsets)

\graphicspath{{Images/}{Charts/}} % Specifies the directory where pictures are stored
\usepackage{times}
\usepackage{pdfpages}
\usepackage{subcaption}
\usepackage[square, numbers, comma, sort&compress]{natbib} % Use the natbib reference package - read up on this to edit the reference style; if you want text (e.g. Smith et al., 2012) for the in-text references (instead of numbers), remove 'numbers'
\usepackage{amsmath}
\usepackage{amssymb}
\usepackage{graphicx}
\usepackage{listings}
\usepackage{siunitx}
\usepackage{float}
\usepackage[euler]{textgreek}
% For algorithms management
\usepackage[linesnumbered, ruled]{algorithm2e}
\SetKwRepeat{Do}{do}{while}%

\lstset{
	language=Python,
	basicstyle=\ttfamily,
	keywordstyle=\color{blue}\ttfamily,
	stringstyle=\color{red}\ttfamily,
	commentstyle=\color{teal}\ttfamily,
	columns=flexible,
	}

\hypersetup{urlcolor=blue, colorlinks=true} % Colors hyperlinks in blue - change to black if annoying

\newcommand\YAMLcolonstyle{\color{red}\mdseries}
\newcommand\YAMLkeystyle{\color{black}\bfseries}
\newcommand\YAMLvaluestyle{\color{blue}\mdseries}

\makeatletter

% here is a macro expanding to the name of the language
% (handy if you decide to change it further down the road)
\newcommand\language@yaml{yaml}

\expandafter\expandafter\expandafter\lstdefinelanguage
\expandafter{\language@yaml}
{
  keywords={true,false,null,y,n},
  keywordstyle=\color{darkgray}\bfseries,
  basicstyle=\YAMLkeystyle,                                 % assuming a key comes first
  sensitive=false,
  comment=[l]{\#},
  morecomment=[s]{/*}{*/},
  commentstyle=\color{purple}\ttfamily,
  stringstyle=\YAMLvaluestyle\ttfamily,
  moredelim=[l][\color{orange}]{\&},
  moredelim=[l][\color{magenta}]{*},
  moredelim=**[il][\YAMLcolonstyle{:}\YAMLvaluestyle]{:},   % switch to value style at :
  morestring=[b]',
  morestring=[b]",
  literate =    {---}{{\ProcessThreeDashes}}3
                {>}{{\textcolor{red}\textgreater}}1
                {|}{{\textcolor{red}\textbar}}1
                {\ -\ }{{\mdseries\ -\ }}3,
}

% switch to key style at EOL
\lst@AddToHook{EveryLine}{\ifx\lst@language\language@yaml\YAMLkeystyle\fi}
\makeatother

\newcommand\ProcessThreeDashes{\llap{\color{cyan}\mdseries-{-}-}}

\title{\ttitle} % Defines the thesis title - don't touch this

\begin{document}
\frontmatter % Use roman page numbering style (i, ii, iii, iv...) for the pre-content pages

\setstretch{1.3} % Line spacing of 1.3

% Define the page headers using the FancyHdr package and set up for one-sided printing
\fancyhead{} % Clears all page headers and footers
\rhead{\thepage} % Sets the right side header to show the page number
\lhead{} % Clears the left side page header

\pagestyle{fancy} % Finally, use the "fancy" page style to implement the FancyHdr headers

\newcommand{\HRule}{\rule{\linewidth}{0.5mm}} % New command to make the lines in the title page
\newcommand{\datasetunimore} {Interactive Museum}
\newcommand{\argmax}{\arg\!\max}
% PDF meta-data
\hypersetup{pdftitle={\ttitle}}
\hypersetup{pdfsubject=\subjectname}
\hypersetup{pdfauthor=\authornames}
\hypersetup{pdfkeywords=\keywordnames}

%----------------------------------------------------------------------------------------
%	TITLE PAGE
%----------------------------------------------------------------------------------------

\begin{titlepage}

\begin{center}
\textsc{\Large \bf \univname}\\[0.5cm] % University name
\textsc{\large Dipartimento di Ingegneria ``Enzo Ferrari"}\\[0.2cm] % Thesis type
%\HRule \\ % Horizontal line
\textsc{\large Corso di Laurea Magistrale in Ingegneria Informatica}\\[5.0cm] % Thesis type

\HRule \\[0.6cm] % Horizontal line
{\huge \bfseries \ttitle}\\[0.4cm] % Thesis title
%\HRule \\[0.6cm] % Horizontal line
%{\huge \bfseries \ttitleITA}\\[0.4cm] % Thesis title
\HRule \\[3.2 cm] % Horizontal line

\begin{minipage}{0.4\textwidth}
\begin{flushleft} \large
\emph{Candidato:}\\
\authornames % Author name - remove the \href bracket to remove the link
\end{flushleft}
\end{minipage}
\begin{minipage}{0.4\textwidth}
\begin{flushright} \large
\emph{Relatore:} \\
\supname \\ % Supervisor name - remove the \href bracket to remove the link
\emph{Correlatore:} \\
\cosupname \\ % Supervisor name - remove the \href bracket to remove the link
\end{flushright}
\end{minipage}\\[5.0cm]

%\large \textit{A thesis submitted in fulfilment of the requirements\\ for the degree of \degreename}\\[0.3cm] % University requirement text
%\textit{in the}\\[0.4cm]
%\groupname\\\deptname\\[2cm] % Research group name and department name

\textsc{\large Anno Accademico 2021-2022}\\[4cm] % Date
%\includegraphics{Logo} % University/department logo - uncomment to place it

\vfill
\end{center}

\end{titlepage}

\clearpage % Start a new page

%----------------------------------------------------------------------------------------
%	SYMBOLS
%----------------------------------------------------------------------------------------

%\clearpage % Start a new page
%
%\lhead{\emph{Symbols}} % Set the left side page header to "Symbols"
%
%\listofnomenclature{lll} % Include a list of Symbols (a three column table)
%{
%$a$ & distance & m \\
%$P$ & power & W (Js$^{-1}$) \\
%% Symbol & Name & Unit \\
%
%& & \\ % Gap to separate the Roman symbols from the Greek
%
%$\omega$ & angular frequency & rads$^{-1}$ \\
%% Symbol & Name & Unit \\
%}
\newcommand{\etal}{\textit{et al.~}}

%----------------------------------------------------------------------------------------
%	ABSTRACT PAGES
%----------------------------------------------------------------------------------------
\addtotoc{Abstract in English} % Add the "Abstract" page entry to the Contents
\englishAbstract{\addtocontents{toc}{\vspace{1em}} % Add a gap in the Contents, for aesthetics
Dental implant placement within the jawbone is a routinely executed surgical
procedure, which can become complex due to the local presence of the Inferior
Alveolar Nerve (IAN). In particular, the nerve is oftentimes in close relation
to the roots of molars, and its position must thus be carefully detailed before
the surgical removal. As avoiding contact with the IAN is a primary concern
during these operations, segmentation plays a key role in surgical
preparations.\\ Unfortunately, 3D manual segmentation requires a huge amount of
time to be carried out, thus perfect anatomical annotations are usually
overlooked in favor of fast executions. The de facto standard in radiology
medical centers for dentistry and maxilla facial purposes is based on sparse
annotations, which can be obtained from 2D images in a relatively small amount
of time. This type of annotation fails to identify a considerable amount of
inner information about the IAN position and the bone structure.
Convolutional Neural Networks are known to provide amazing results in the
segmentation task and could provide to medical operators an automatic and
precise way to detect the IAN in patients, allow them to be more accurate during
the operations where this information is critical. This type of network requires
a large amount of labeled data to reach a good level of accuracy while in the
medical field is hard and expensive to acquire, meanwhile a good amount of
unlabeled data is available.\\
In this work we study and analyze the effectiveness of some of the most recent
proposal made in the literature for segmentation in the medical imaging field,
some of which have recently obtained state-of-the-art results for other similar
tasks. Even if we do not reach the same level of accuracy we expected, we show
that we were able to improve the speed of the training process and the
performance of the network, while keeping the same level of accuracy of the
original network.

%Depositate anche su esse3, quindi se le cambi cambia anche lì
\textbf{Keywords:} Medical Imaging, Deep Learning, 3D Segmentation, Computer Vision, Cone Beam Computed Tomography.
}
\clearpage % Start a new page

\addtotoc{Abstract in Italian} % Add the "Abstract" page entry to the Contents
\italianAbstract{\addtocontents{toc}{\vspace{1em}} % Add a gap in the Contents, for aesthetics
L'inserimento di impianti dentali nella mascella è una ruotine chirurgica
abbastanza comune, ma è resa non banale dalla presenza del nervo alveolare
inferiore. Questo nervo è posizionato in prossimità delle radici dei molari,
e la sua posizione deve essere accuratamente delineata prima di poter eseguire
l'operazione voluta, in quanto danneggiare tale nervo potrebbe essere causa di
gravi conseguenze per il paziente, come ad esempio la perdita della
sensibilità o la paralisi di tale zona. Poichè evitare il contatto con questo
nervo è essenziale, un'accurata segmentazione di tale nervo è un requisito
chiave. Sfortunatamente, una segmentazione 3D manuale richiede un'enorme
quantità di tempo, per questo motivo solitamente si preferiscono metodi meno
precisi ma più veloci e per questo motivo lo standard de-facto per i dentisti è
basato sull'effettuare una annotazione sparsa, ottenuta da immagini 2D. Questo
tipo di annotazione fallisce nell'identificare una buona parte delle
informazioni che riguardano la posizione del nervo e la struttura ossea del
canale in cui esso è contenuto.\\
Le reti neurali convolutive hanno dimostrato di poter ottenere risultati ottimi
nell'effettuare la segmentazione di immagini e in questo caso potrebbe offrire
ai medici un metodo automatico, preciso e veloce di delineare il nervo alveolare
inferiore, essenziale soprattutto durante le operazioni più critiche. Questo
tipo di reti neurali richiede una quantità enorme di dati annotati per poter
imparare a svolgere correttamente questo tipo di task con una accuratezza
soddisfacente. In ambito medico, è possibile ottenere una decente quantità di
immagini, ma non di annotazioni. Per questo motivo è necessario progettare
accuratamente i metodi utilizzati durante il design e l'allenamento di tali reti
in modo da superare questa mancanza senza perdere in prestazioni.
In questa
tesi, diverse tecniche molto recenti vengono provate ed analizzate, per
studiarne l'efficacia in questo specifico lavoro. Alcuni metodi che recentemente
hanno avuto molto successo nell'ambito della computer vision, non hanno avuto
gli stessi risultati sperati in questo specifico caso, mentre altri approcci si
sono rivelati molto efficaci nell'ottimizzare e velocizzare il processo di
training della rete, senza perdere accuratezza nella segmentazione.

\textbf{Parole Chiave:} Medical Imaging, Deep Learning, 3D Segmentation, Computer Vision, Cone Beam Computed Tomography.
}
\clearpage % Start a new page

%----------------------------------------------------------------------------------------
%	SINTESI IN ITALIANO PAGES
%----------------------------------------------------------------------------------------

\addtotoc{Summary in Italian} % Add the "Abstract" page entry to the Contents
%
\italianSummary{\addtocontents{toc}{\vspace{1em}} % Add a gap in the Contents, for aesthetics
%
Il regolamento della Facoltà di Ingegneria di Modena prevede che le tesi scritte
in lingua inglese debbano contenere un \textit{abstract} ed un'ampia sintesi dei
contenuti in lingua italiana: in accordo a questa regola, proponiamo un sunto
degli argomenti, delle tecniche e dei risultati che verranno delineati
nell'elaborato. Si noti che non si tratta di un sunto esaustivo, e che non è
possibile valutare la tesi dalla semplice lettura di queste righe. Per una
descrizione più dettagliata e più rigorosa, e per i risultati sperimentali
ottenuti, si rimanda al testo in inglese.

Questa tesi tratta della segmentazione del canale alveolare inferiore su
immagini di tomografia conica a fascio di coni (Cone Beam Computed Tomography,
CBCT) tramite l'utilizzo di reti neurali.\\
}
%
\clearpage % Start a new page

%----------------------------------------------------------------------------------------
%	DEDICATION
%----------------------------------------------------------------------------------------
\mainmatter
\setstretch{1.3} % Return the line spacing back to 1.3

\pagestyle{empty} % Page style needs to be empty for this page

\dedicatory{To ...} % Dedication text

\addtocontents{toc}{\vspace{2em}} % Add a gap in the Contents, for aesthetics

%----------------------------------------------------------------------------------------
%	ACKNOWLEDGEMENTS
%----------------------------------------------------------------------------------------

\setstretch{1.3} % Reset the line-spacing to 1.3 for body text (if it has changed)

\acknowledgements{\addtocontents{toc}{\vspace{1em}} % Add a gap in the Contents, for aesthetics

Foremost, I would like to express my sincere gratitude to ...
}
\clearpage % Start a new page

%----------------------------------------------------------------------------------------
%	LIST OF CONTENTS/FIGURES/TABLES PAGES
%----------------------------------------------------------------------------------------

\pagestyle{fancy} % The page style headers have been "empty" all this time, now use the "fancy" headers as defined before to bring them back

\lhead{\emph{Contents}} % Set the left side page header to "Contents"
\tableofcontents % Write out the Table of Contents

\lhead{\emph{List of Figures}} % Set the left side page header to "List of Figures"
\listoffigures % Write out the List of Figures

\lhead{\emph{List of Tables}} % Set the left side page header to "List of Tables"
\listoftables % Write out the List of Tables

\lhead{\emph{List of Listings}} % Set the left side page header to "List of Listings"
\lstlistoflistings % Write out the List of Listings

%----------------------------------------------------------------------------------------
%	ABBREVIATIONS
%----------------------------------------------------------------------------------------

%\clearpage % Start a new page
%
%\setstretch{1.5} % Set the line spacing to 1.5, this makes the following tables easier to read
%
%\lhead{\emph{Abbreviations}} % Set the left side page header to "Abbreviations"
%\listofsymbols{ll} % Include a list of Abbreviations (a table of two columns)
%{
%\textbf{LAH} & \textbf{L}ist \textbf{A}bbreviations \textbf{H}ere \\
%%\textbf{Acronym} & \textbf{W}hat (it) \textbf{S}tands \textbf{F}or \\
%}

%----------------------------------------------------------------------------------------
%	PHYSICAL CONSTANTS/OTHER DEFINITIONS
%----------------------------------------------------------------------------------------

%\clearpage % Start a new page
%
%\lhead{\emph{Physical Constants}} % Set the left side page header to "Physical Constants"
%
%\listofconstants{lrcl} % Include a list of Physical Constants (a four column table)
%{
%Speed of Light & $c$ & $=$ & $2.997\ 924\ 58\times10^{8}\ \mbox{ms}^{-\mbox{s}}$ (exact)\\
%% Constant Name & Symbol & = & Constant Value (with units) \\
%}


%----------------------------------------------------------------------------------------
%	THESIS CONTENT - CHAPTERS
%----------------------------------------------------------------------------------------

 % Begin numeric (1,2,3...) page numbering

\pagestyle{fancy} % Return the page headers back to the "fancy" style

% Include the chapters of the thesis as separate files from the Chapters folder
% Uncomment the lines as you write the chapters

% Chapter 1

\chapter{Segmentation of the Inferior Alveolar Canal: an overview}

\lhead{Chapter 1. \emph{Segmentation of the Inferior Alveolar Canal}} % This is for the header on each page - perhaps a shortened title
\label{chp:introductive_chapter}
%----------------------------------------------------------------------------------------

\def\:{\hskip0pt} %Definisce un modo veloce per permettere a latex di sillabare correttamente anche parole come 4-connectivity. Il corretto utilizzo è il seguente: 4\:-\:connectivity.

\section{Introduction}
Dental implant placement within the jawbone is a routine surgical procedure that
can become complicated due to the presence of the Inferior Alveolar Nerve (IAN)
nearby. The nerve, in particular, is frequently in close proximity to the roots
of molars, and its position must thus be meticulously detailed prior to surgical
removal. Avoiding contact with the IAN is a primary concern during these
operations, thus its segmentation is crucial in surgical planning.
Today the standard de-facto is to take a CBCT scan of the jawbone and a 2D
panomaric view is extracted. This view allow medical experts to depict the IANs
position with line. We refer to this type of annotation as \emph{sparse} or
\emph{2D} annotation. A 3D annotation of the IAC is often avoided as it would
require a huge amount of time, but this type of segmentation would offer a much
precise knowledge of the position of the IAN and IAC and could allow dentists to
plan a more detailed surgical approach. For this reason a lot of research about
automatic segmentation of the IAC has been carried out and is still active today.\\

In this chapter we first describe in details the role of the IAN and the IAC,
what a CBCT is and the definitions and carachteristics of different types of
segmentations.
In the following chapter we will take a look on how segmentation is performed
today using neural networks paying more attention field of medical images.
Next in chapter 3 we will detail the dataset, the network, the metrics, and
other tools used as baseline to define the current state of the art, which is
essential to understand the importance and goodness of the work that i've
carried out. Last, in chapter 4, all the different techniques tried will be
described in details, to conclude with a description of which would possible be
the future work for this specific tasks.

\section{Inferior Alveolar Canal}
The Inferior Alveolar Canal (IAC) is a small passageway shaped as a tube that
runs through the lower jawbone. It houses the Inferior Alveolar Nerve (IAN),
which is responsible for transmitting sensory information from the teeth, gums
and lips to the brain. It also provides motor innervation to the muscles of
mastication (i.e. the muscles responsible for chewing). Dentists need to be able
to accurately locate the IAC before performing certain surgical operations, such
as tooth extractions or placement of dental implants. This is because the IAC is
located very close to the roots of the teeth, and damage to the IAC during
surgery can result in permanent nerve damage which would cause numbness,
tingling, and pain in the affected area. In severe cases, it can also lead to
muscle weakness and paralysis.

\section{Cone Beam Computed Tomography}
Cone beam computed tomography (CBCT) is a medical imaging technique consisting
of X-ray computed tomography where rays are divergent, forming a cone. This type
of computed tomography is well suited for imaging the craniofacial area as it
provides clear images of highly contrasted structures, very helpfull to
evaluate bones. It has become common in dentistry such as oral surgery,
endodontics and orthodontics.
The main reasons and advantages of CBCT with respect to other CTs are:
\begin{enumerate}
  \item{\textbf{X-ray beam limitation:} reducing the size of the irradiated area
  by collimating the primary x-ray beam to the are of interest minimizes the
  radiation dose. Most CBCT units can be adjusted to scan small regions for
  specific diagnostic task. They are also able to scan the whole craniofacial
  structure if needed.}
  \item{\textbf{Image accuracy:} We created a novel, large, and publicly
  available maxillo-facial CBCT (Cone Beam Computed Tomography) dataset, with 2D
  and 3D manual annotations, provided by expert clinicians. All CBCT units
  provide voxel resolutions that are isotropic (i.e. equals in all the 3
  dimensions) while in conventional CT, voxel are anisotripic (i.e. rectangular
  cubes).}
  \item{\textbf{Rapid scan time:} CBCT acquires all the basis images in a single
  rotation, thus scan time goes from 10s to 70s. Although faster scanning time
  usually means fewer basis images from which to reconstruct the volumetric
  dataset, motion artifacts due to subject movement are reduced.}
\end{enumerate}
These advantages come with some drawbacks: Hounsfield units (HU) is the metric
used to determine the radiodensity of tissue analized. In the Hounsfield scale,
numbers go from values of $-1000$ for air to values of $1600$ for dense bones.
In CBCT scans, the radiodensity is inaccurate because different areas in the
scan appear with different greyscale values depending on their relative
positions in the organ being scanned, despite possessing identical densities,
because the image value of a voxel of an organ depends on the position in the
image volume. HU measured from the same anatomical area with both CBCT and
medical-grade CT scanners are not identical and are thus unreliable for
determination of site-specific, radiographically-identified bone density for
purposes such as the placement of dental implants, as there is "no good data to
relate the CBCT HU values to bone quality" \cite{Miles2007}.\\
The images resulting from a CBCT scans are usually exported as DICOM (Digital
Imaging and Communications in Medicine) which is the standard used worldwide to
store, exchange, and transmit medical images.

\section{DICOM file format}
TODO?

\section{Image Segmentation}
Image segmentation is a well known topic in computer and image processing with a
wide range of application, such as medical imaging, robotics, video
surveillance, etc.
It involves partitioning images into one or more objects and can also includes
classify these objects. Many traditional algorithms have been developed in the
literature but, in the most recent years, they have all been dominated by deep
neural networks. Since the 2015 a huge amount of different types of networks
that aim to perform image segmentation has been proposed for each of the field
where it's needed.
Before presenting how nowday segmentation is performed, we must state which are
the different type of segmentation that have been classified:

\begin{itemize}
  \item{\textbf{Semantic segmentation:}
  Semantic Segmentation perform a pixel-by-pixel classification with a predefined
  set of objects categories for all the pixels of the images. In pratice, given a
  RGB image (\texttt{height} $\times$ \texttt{width} $\times$ \texttt{3}) we output
  a segmentation map of size (\texttt{height} $\times$ \texttt{width} $\times$
  \texttt{classes}) where each value correspond to which class the same pixel in the
  original images belongs.}

  \item{\textbf{Instance segmentation:}
  One possible issue of semantic segmentation is that it doesn't allow to
  distinguish two or more object of the same class when they overlap in the image.
  Instance segmentation overcome this problem by outputting a different number of
  channels based on the number of instances present in the image.}

  \item{\textbf{Panoptic segmentation:}
  The latter type of segmentation is called Panoptic segmentation and is the
  result of the previously presented method joined together. The difference with
  the instance segmentation is that in this case instances are not allowed to
  overlap then for to a single pixel, a single instance must be assigned.}
\end{itemize}

% Figure 1: example of a segmentation on mri image
\begin{figure}[h]
  \centering
  \includegraphics[width=0.8\textwidth]{Images/example_segmentation.png}
  \caption{Example of a multiclass semantic segmentation, different colors represent different classes}
  \label{fig:segmentation}
\end{figure}



% \section{Segmentation using Deep Neural Networks}
% As stated before, nowdays deep learning is the dominant method used for
% segmentation. Convolutional Neural Networks are still the most used type of
% networks used to date but new promising technique, as Graph Neural Networks or
% Transformers, are always more and more used.

% Chapter 2

\chapter{Segmentation Neural Networks}

\lhead{Chapter 2. \emph{Segmentation Neural Networks}} % This is for the header on each page - perhaps a shortened title
\label{chp:cbct}
%----------------------------------------------------------------------------------------

\def\:{\hskip0pt} %Definisce un modo veloce per permettere a latex di sillabare correttamente anche parole come 4-connectivity. Il corretto utilizzo è il seguente: 4\:-\:connectivity.

\section{Segmentation using Deep Neural Networks}
\label{sec:segmentation}
Today, deep neural netoworks are the state-of-the-art in many fields, including
image segmentation. In this section, we will briefly review the most common
approaches to segmentation using deep neural networks, and we will discuss the
advantages and disadvantages of each approach.

\subsection{Fully Convolutional Networks}
\label{sec:fcn}
Fully convolutional networks (FCNs) are a class of deep neural networks that
are designed to perform pixel-wise classification. The main idea behind FCNs is
to use a convolutional neural network (CNN) to extract features from the input
image and return an output of the same size as the input image, where each
pixel is assigned a class label. The main advantage of FCNs is that they can be
trained end-to-end, which means that the network can be trained to perform the
classification of each pixel without the need of any post-processing step. On
the other hand, Deep Neural Networks lacks for explainability, which is a major
drawback for medical applications.
The architecture of such networks can be grouped as shown in Figure \ref{fig:fcn_architecture}.
\begin{figure}[ht]
  \centering
  \includegraphics[width=0.8\linewidth]{Images/seg-arch.png}
  \caption{Groups of a segmentation network architecture.}
  \label{fig:fcn_architecture}
\end{figure}

\subsection{Convolutional Neural Networks}
A convolutional neural network or CNN consists of a stack of
three main neural layers: convolutional layer, pooling layer, and fully
connected layer. Each layer has its own role. The convolution layer
detects distinct features like edges or other visual elements in an image.
Convolution layer performs mathematical operation of multiplication of local
neighbours of an image pixel with kernels. CNN uses different kernels for
convolving the given image for generating its feature maps. Pooling layer
reduces the spatial \texttt{(width, height)} dimensions of the input data for the next
layers of neural network. It does not change the depth of the data. This
operation is called as subsampling. This size reduction decreases the
computational requirements for upcoming layers. The fully connected layers
perform high-level reasoning in NN. These layers integrate the various feature
responses from the given input image so as to provide the final results.\\
Different CNN models have been reported in the literature, including AlexNet,
GoogleNet, VGG, Inception, SequeezeNet, and DenseNet. This type of networks were
mostly used for classification, but they have been easily adapted to perform
segmentation.

\subsection{Fully Convolutional Networks}
In fully convolutional network (FCN), only convolutional layers exist. The
different existing in CNN architectures can be modified into FCN by converting
the last fully connected layer of CNN into a fully convolutional layer. This
type of networks can output spatial segmentation map and can have dense
pixel-wise prediction from the input image of full size instead of performing
patch-wise predictions. They can also uses skip connections, when performing upsampling
on feature maps from final layer, these skip connections fuses it with the
feature map of previous layers. The model thus produces a detailed segmentation
in just one go but as drawback they do not have a global context of the image
and the output can be fuzzy close to the boundaries of segmented objects.

\subsection{Encoder-Decoder Networks}
Encoder-decoder based models employ two-stage model to map data points from the
input domain to the output domain. The encoder stage compresses the given input
to a latent space representation, while the decoder predicts the output from
this representation. This latent space representation is a compressed version of
the input image, which is used to generate the output, it have smaller spatial dimension
than the input image but an increased depth. In order to upsample this latent space representation
to the size of the input image, transposed convolutional layers are used.
One of the most popular encoder-decoder networks is the U-Net
\cite{ronneberger2015u}, for which we can find in literature many variants. The
U-Net architecture is shown in Figure \ref{fig:unet_architecture}.
% \begin{figure}[ht]
%   \centering
  % \includegraphics[width=0.8\linewidth]{Images/unet-arch.png}
  % \caption{U-Net architecture.}
  % \label{fig:unet_architecture}
% \end{figure}
U-Net model has a downsampling and upsampling part. The downsampling
section with FCN like architecture extracts features using $3 \times 3$ convolutions to
capture context. The upsampling part performs deconvolution to decrease the
number of computed feature maps. The feature maps generated by downsampling or
contracting part are fed as input to upsampling part so as to avoid any loss of
information. The symmetric upsampling part provides precise localization. The
model generates a segmentation map which categorizes each pixel present in the
image. This type of architecture proposed in 2015 is still widely used today as
it can obtain state-of-the-art resuts.

\subsection{Regional Convolutional Neural Networks}
Regional convolutional network has been utilized for object detection and
segmentation task. The R-CNN architecture presented in [69] generates region
proposal network for bounding boxes using selective search process. These region
proposals are then warped to standard squares and are forwarded to a CNN so as
to generate feature vector map as output. The output dense layer consists of
features extracted from the image and these features are then fed to
classification algorithm so as to classify the objects lying within the region
proposal network. The algorithm also predicts the offset values for increasing
the precision level of the region proposal or bounding box. The processes
performed in R-CNN architecture are shown in Figure 4. The use of basic RCN
model is restricted due to the following:

\subsection{DeepLab}
DeepLab model employs pretrained CNN model ResNet-101/VGG-16 with atrous
convolution to extract the features from an image. The use of atrous
convolutions gives the following benefits:
\begin{itemize}
  \item{It controls the resolution of feature responses in CNNs}
  \item{It converts image classification network into a dense feature extractor
    without the requirement of learning of any more parameters employs
    conditional random field (CRF) to produce fine segmented output}
\end{itemize}
The various variants of DeepLab have been proposed in the literature including
DeepLabv1, DeepLabv2, DeepLabv3, and DeepLabv3+.
In DeepLabv1, the input image is passed through deep CNN layer with one or
two atrous convolution layers. This generates a coarse feature map. The feature
map is then upsampled to the size of original image by using bilinear
interpolation process. The interpolated data is applied to fully connect
conditional random field to obtain the final segmented image.

\section{Dealing with 3D Medical Images}
It is not uncommon in the medical field to deal with 3D images that comes from
CT scans or MRI scans. In this case, a 3D image can be represented as a stack of
2D images and fed them to the network. Each output is then stacked back to
produce the final 3D output. Another approach is to use 3D convolutional
networks. The 3D convolutional network is a generalization of the 2D
convolutional network to 3D data.
In the first case, each 2D layer is indipendend from the others, while in the
second case, the 3D convolutional network is able to learn the spatial
relationships between the different slices of the 3D image.


% Chapter 3

\chapter{Image Segmentation}

\lhead{Chapter 3. \emph{Image Segmentation}} % This is for the header on each page - perhaps a shortened title
\label{chp:segmentation}
%----------------------------------------------------------------------------------------

\def\:{\hskip0pt} %Definisce un modo veloce per permettere a latex di sillabare correttamente anche parole come 4-connectivity. Il corretto utilizzo è il seguente: 4\:-\:connectivity.

\section{Introduction}
Image segmentation is a well known topic in computer and image processing with a
wide range of application, such as medical imaging, robotics, video
surveillance, etc.
It involves partitioning images into one or more objects and can also includes
classify these objects. Many traditional algorithms have been developed in the
literature but, in the most recent years, they have all been dominated by deep
neural networks. Since the 2015 a huge amount of different types of networks
that aim to perform image segmentation has been proposed for each of the field
where it's needed.
Before presenting how nowday segmentation is performed, we must state which are
the different type of segmentation that have been classified.

\subsection{Semantic Segmentation}
Semantic Segmentation perform a pixel-by-pixel classification with a predefined
set of objects categories for all the pixels of the images. In pratice, given a
RGB image (\texttt{height} $\times$ \texttt{width} $times$ \texttt{3}) we output
a segmentation map of size (\texttt{height} $\times$ \texttt{width} $\times$
\texttt{1}) where each value correspond to which class the same pixel in the
original images belongs.

\subsection{Instance segmentation}
One possible issue of semantic segmentation is that it doesn't allow to
distinguish two or more object of the same class when they overlap in the image.
Instance segmentation overcome this problem by outputting a different number of
channels based on the number of instances present in the image.

\subsection{Panoptic segmentation}
The latter type of segmentation is called Panoptic segmentation and is the
result of the previously presented method joined together. The difference with
the instance segmentation is that in this case instances are not allowed to
overlap then for to a single pixel, a single instance must be assigned.

\subsection{Segmentation using Deep Neural Networks}
As stated before, nowdays deep learning is the dominant method used for
segmentation. Convolutional Neural Networks are still the most used type of
networks used to date but new promising technique, as Graph Neural Networks or
Transformers, are always more and more used.

% Chapter 4

\chapter{Code Refactoring}

\lhead{Chapter 4. \emph{Code Refactoring}} % This is for the header on each page - perhaps a shortened title
\label{chp:refwork}
%----------------------------------------------------------------------------------------

\def\:{\hskip0pt} %Definisce un modo veloce per permettere a latex di sillabare correttamente anche parole come 4-connectivity. Il corretto utilizzo è il seguente: 4\:-\:connectivity.
\section{Introduction}
Cipriano et al. spent some effort developing the dataset and a to design a
network which was capable to perform the segmentation of the Inferior Alveolar
Canal.
Together with the novel dataset, they proposed a modified version of
U-Net 3D as a backbone for both the label propagation and the IAC segmentation.
Both these networks and methods were introduced in the previous
chapters and will be explained more in detail in the following sections.
As the code was written near the deadline for the submission of the paper,
the authors did not have focused on producing a codebase that was optimized for
reusability, extensibility, and efficiency.
For this reason, my work starts by refactoring the codebase to add these features
which improved all my successive experiments.

\section{Reference Work}
The reference work that has been taken as starting point for this work has been
made by Cipriano et al. in 2022 \cite{cipriano2022improving}.
The pipeline proposed can be divided into two main steps: the deep label
propagation and the IAC segmentation.

\subsection{Positional PadUNet3D}
A slightly modified version of 3D U-Net has been proposed as a backbone for both
steps. Differently from the original 3D U-Net architecture, every
three-dimensional convolution in our CNN applies 2 pixels padding along each
dimension. Although this alteration does not cause any variation in terms of
performance, it ensures that the output of each convolution has the same size as
its input along axis \emph{x}, \emph{y}, and \emph{z}.

Resolution changes are therefore due uniquely to the three max-pooling layers,
each halving the size of the volumes. When using the aforementioned input size,
the encoder output is composed of $512$ feature maps of size $10 \times 10
\times 10$. Inside the decoder, on the other hand, resolution changes are caused
by transposed convolutions, with dimensions $2 \times 2 \times 2$ for both the
kernel and stride to double the size of the feature maps. This adjustment
ensures resolution symmetry between features in the decoder and corresponding
maps in the encoder, thus allowing them to be simply concatenated with skip
connections. Moreover, the output of the model naturally has the same dimensions
as the input. The final output is a single channel volume, to which we apply a
Sigmoid activation function and a threshold at $0.5$ to obtain the final binary
prediction mask.

The segmentation architecture is further enriched with a positional embedding.
Since our sub-volumes are extracted from the original scan following a fixed
grid, we exploit positional information derived from the location of the
top-left and bottom-right corners of the sub-volume. Specifically, these global
coordinates are fed to a linear layer which yields a single feature map of
dimensions $10 \times 10 \times 10$. This positional embedding gets concatenated
to the output of the encoder, and then fed to the decoder. Exploiting positional
information ensures two extremely important benefits:
\begin{itemize}
  \item{During training, the CNN is fed with implicit information about areas
  close to the edges of the scan, where the IAN is very unlikely to be present.
  This piece of knowledge greatly reduces the number of false positives during
  inference. As a matter of fact, no postprocessing method is required to refine
  the output of the proposed CNN}
  \item{Information about cut positions helps the network to better shape the
  output: sub-volumes located close to the mental foramen generally present a
  much thinner canal than those located in the mandibular foramen. This
  technique could indeed be employed for several classes of medical data, as
  anatomical structures can substantially vary according to their location.}
\end{itemize}

\subsection{Deep Label Propagation}
The presence of a new 3D voxel-level annotated dataset paves the way to a new
frontier for label propagation: we can indeed employ 3D annotations to supervise
a deep label propagation neural network, trained to expand sparse labels into
dense ones, and produce high-quality synthetic ground truths for the
segmentation task. We christen this new label propagation approach Deep Label
Expansion, or Deep Expansion in short. The deep expansion model is based on the
proposed segmentation network, Positional PadUNet. The main difference regards
the input layer, which is changed in order to accept a concatenation of both the
raw volume data and the sparse annotations, rendered as a binary channel. Thus,
the network input has a shape of $2 \times 80 \times 80 \times 80$. Same as
Positional PadUNet, a positional embedding is concatenated to the encoder
output. By means of this model, we generate a synthetic dataset which will be
employed to pre-train our final segmentation CNN. Once again, the positional
embedding supplies important information about the location of the cut, which is
closely related to the diameter of the expanded labeled canal.

\section{Refactoring}
The code was divided into two main repositories, one for the label propagation
task and one for the IAC segmentation task. For this reason, a lot of code was
duplicated and a single repository was created to host the whole codebase.
Moreover, the code has been rewritten from scratch, and only some parts of the
original codebase have been copied.
Heavy use of the TorchIO library, mantained by Perez Garcia
\cite{PerezGarcia2021torchio}, has been made to simplify and optimize part of
the pipeline.
% restructured to allow the user to easily change the parameters of the network

\subsection{Data Loading}
The first issue tackled was the lack of use of TorchIO classes as an abstraction
for the new proposed dataset. The original codebase was using a custom class that
was responsible to load the data, compute some statistics, and reading part of
the config file to select which transformation to apply to perform the data
augmentation. Moreover, custom transformation functions were defined here in the
same file.

One of the main advantages of using TorchIO is that it provides three main
data structures which are used to represent the data: \emph{Image},
\emph{Subject} and \emph{SubjectsDataset}.

The Image class, representing one medical image, stores a 4D tensor, whose
voxels encode, e.g., signal intensity or segmentation labels, and the
corresponding affine transform, typically a rigid (Euclidean) transform, to
convert voxel indices to world coordinates in mm. Arbitrary fields such as
acquisition parameters may also be stored. Subclasses are used to indicate
specific types of images, such as \emph{ScalarImage} and \emph{LabelMap}, which
are used to store, e.g., CBCT scans and segmentations, respectively.

The Subject is a data structure used to store images associated with a subject
and any other metadata necessary for processing. It is essential to link a
ScalarImage to its relative LabelMaps.

The SubjectsDataset is a subclass of torch.utils.data.Dataset which is used to
store a collection of Subjects. It can be used with PyTorch DataLoader for
efficient loading and augmentation. The figure \ref{fig:subjectsdataset}
depicts the relationship between these classes.
\begin{figure}[h]
  \centering
  \includegraphics[width=0.9\linewidth]{Images/subjectsdataset.png}
  \caption{SubjectsDataset class and its relationships.}
  \label{fig:subjectsdataset}
\end{figure}

With these data structures, a Maxillo class, which inherit from
SubjectsDataset has been created to load the data. This class takes as input the
path of the dataset, which is split to load (train, validation, test, synthetic)
and optionally a list of Transformations that must be used as data augmentations. This class is completely independent from the rest of the code, and thus can be used in any other project which requires loading the Maxillo dataset.

\subsection{Data Augmentation}
Transformation functions that were previously defined in the same file of the
data loader have been moved to a new file. Some of them were already implemented
inside TorchIO and thus have been removed. The remaining ones have been adapted
to inherit from the abstract \texttt{class torchio.transforms.Transform}, in
order to be easily used inside TorchIO, such with the \emph{Composability}
method which allows creating a directed acyclic graph of transformations that
must be applied to the data.

\subsection{Config and command line arguments}
In the original codebase, the parameters of the program were passed both as
arguments from the command line and in the config file. This could generate some
confusion and make it harder to reproduce experiments as the config file alone
wasn't enough. For this reason, the config file is now the only source of
parameters, while the only arguments that the software accepts are the path of
the config file and a boolean value to enable/disable Tensorboard \cite{abadi2015tensorflow}.

% An example of a config file is shown below:

% \begin{lstlisting}[language=YAML,breaklines=true,caption=Example of config file., label=union]
% title: 'canal_generator_train'
% project_dir: '/homes/llumetti/results'
% seed: 47
% tensorboard_dir: '/homes/llumetti/tensorboard'

% experiment:
%   name: 'Generation'

% data_loader:
%   dataset: '/nas/softechict-nas-1/llumetti/maxillo'
%   augmentations: 'configs/augmentations.yaml'
%   background_suppression: 0
%   batch_size: 2
%   labels:
%     BACKGROUND: 0
%     INSIDE: 1
%   mean: 0.08435
%   num_workers: 8
%   patch_shape:
%   - 120
%   - 120
%   - 120
%   resize_shape:
%   - 168
%   - 280
%   - 360
%   sampler_type: grid
%   grid_overlap: 0
%   std: 0.17885
%   volumes_max: 2100
%   volumes_min: 0
%   weights:
%   - 0.000703
%   - 0.999

% model:
%   name: 'PosPadUNet3D'

% loss:
%   name: 'Jaccard'

% lr_scheduler:
%   name: 'Plateau'

% optimizer:
%   learning_rate: 0.1
%   name: 'SGD'

% trainer:
%   reload: False
%   checkpoint: null
%   do_train: False
%   do_test: False
%   do_inference: True
%   epochs: 100
% \end{lstlisting}

\subsection{Interfaces}
Something about interfaces used for Models, Losses, Optimizers, Schedulers and
Experiments...

\subsection{Deploying and Running}
The code was managed using Git and the repository is hosted on GitHub the
repository is public and can be found at
\url{https://github.com/AImageLab-zip/alveolar\_canal}. By using Git, it was
possible to easily deploy the code on the AImagelab Server by means of a simple
script that keeps the code up to date with the main branch. Other branches were
used to develop new features and to test them before merging everything in the
main branch. The Gitflow workflow was used to manage the branches. The training
was scheduled using the SLURM scheduler \cite{yoo2003slurm}, which is installed
on the server. Some scripts were written to automatize the training process,
which often required running sequentially multiple experiments with different
configuration files and to recover training from a given checkpoint when the
training was long and the scheduler kill the process for exceeding the time
limit.


% Chapter 5

\chapter{Trying to improve the State of the Art}

\lhead{Chapter 5. \emph{Trying to improve the State of the Art}} % This is for the header on each page - perhaps a shortened title
\label{chp:refwork}
%----------------------------------------------------------------------------------------

\def\:{\hskip0pt} %Definisce un modo veloce per permettere a latex di sillabare correttamente anche parole come 4-connectivity. Il corretto utilizzo è il seguente: 4\:-\:connectivity.
The architecture proposed by Cipriano et al. was quite simple yet effective. The
main novel addition was the use of positional encoding. We now want to try to
improve the results obtained by them by adding some more complex
and more recent architectures, losses, and optimizers.

\section{Boundary Loss}
Jointly with PosPadUNet3D, standard DICE loss was used. We also tried to use
CrossEntropy loss, but the results were not as good as with DICE loss so it has
been discarded. As we were dealing with highly unbalanced data, less than 1\% of
the whole volume is the canal, we looked for alternative losses which deal with
this problem.

Kervadec et al. \cite{kervadec2019boundary} proposed a loss function that takes
the form of a distance metric on the space of contours instead of regions. Dice
or cross-entropy, are based on regional integrals, which are convenient for
training deep neural networks. In practice, these regional integrals are
summations over the segmentation regions of differentiable functions, each
directly invoking the softmax probability outputs of the network. Therefore,
standard stochastic optimizers such as SGD are directly applicable.
Unfortunately, difficulties occur for highly unbalanced segmentations, for
instance, when the size of the target foreground region is several orders of
magnitude less than the background size, which is a common characteristic for
medical images. The problem with such losses is that they assume identical
importance distribution for all the samples and classes.

In the aforementioned paper, the authors proposed a new type of loss, named
\emph{Boundary loss} that aims to mitigate the issues related to regional losses
in highly unbalanced segmentation problems. Rather than using unbalanced
integrals over the regions, a boundary loss uses integrals over the boundary
between the regions. Furthermore, it provides information that is complementary
to regional losses. It is, however, challenging to represent the boundary points
corresponding to the regional softmax outputs of a CNN. This difficulty may
explain why boundary losses have been avoided in the context of deep
segmentation networks.

\subsection{Formulation}
The Boundary loss has been formulated as follows: let $I: \Omega \subset
\mathbb{R}^{2,3} \rightarrow \mathbb{R}$ denotes an image with spatial domain
$\Omega$, and $g: \Omega \rightarrow {0,1}$ a binary ground truth segmentation
of the image such that $g(p) = 1$ if the pixel $p$ belongs to the target region
$G \subset \Omega$ and $0$ otherwise. Let $s_\theta : \Omega \rightarrow [0,1]$
denotes the softmax probability output of a deep segmentation network, and
$S_\theta \subset \Omega$ denotes the corresponding segmentation region:
$\S_\theta = {p \in \Omega | s_\theta(p) >= \delta}$ for some threshold
$\delta$. Let $\Delta G$ denote a representation of the boundary of the
ground-truth region $G$ (i.e. the set of points of $G$, which have a spatial
neighbor in background $\Omega \setminus G$) and $\Delta S_\theta$ denoting the
boundary of the segmentation region defined by the network output.\\
The boundary loss can now be defined as:

\begin{equation}
  \label{eq:boundaryloss}
  \mathcal{L}_{\text{boundary}}(\theta) = \int_{\Omega} \phi_G(q)s_\theta(q)dq
\end{equation}

Where $\phi_G$ is the level set of the ground-truth region $G$, obtained by
using the signed distance transform over $G$.

\subsection{Combining with other loss functions}
As we have already said, this type of loss is complementary to the regional
losses, since it provides information about the boundary of the segmentation
region, therefore they can be combined. In our experiments, we used the
following combination:
\begin{equation}
  \label{eq:boundaryloss}
  \mathcal{L}(\theta) = (1-\alpha)\mathcal{L}_{\text{boundary}}(\theta) +
  \alpha\mathcal{L}_{\text{DICE}}(\theta)
\end{equation}
Where $\alpha$ is a hyperparameter that weights the two losses.

Kervadec et al. \cite{kervadec2019boundary} in the paper where they presented
this type of novel loss they've adopted a peculiar way to set this $\alpha$. As
the training was kind of unstable they set $\alpha$ to an initial value of $0$
and increased its value by $0.01$ and each training epoch until a fixed
threshold. This type of approach have given the best results in terms of
stability and performance thus we decided to adopt it also in our experiments.

\section{PosDeepLab v3+}
DeepLab v3+ is another well-known architecture for semantic segmentation. It was
proposed by Chen et al. in 2015 with version 1 \cite{chen2014deeplab1}. Later in
the years, they proposed small changes to the architecture, which resulted in
version 2 \cite{chen2017deeplab2}, 3 \cite{chen2017deeplab3}, and the latest
one, version 3+ \cite{chen2018deeplab3p}. The main novelty that can be found in
this architecture is the use of Atrous Spatial Pyramid Convolutions, which are
dilated convolutions in parallel, each one with a different dilation rate and
later concatenated, the use of a fully connected Conditional Random Field, which
improves the segmentation by maximizing the label agreement between similar
pixels, and the use of depth-wise separable convolutions, which are a way to
reduce the number of parameters in the network. The model architecture is shown
in Figure \ref{fig:deeplabv3+}.
\begin{figure}[h]
  \centering
  \includegraphics[width=0.8\textwidth]{Images/DeepLabv3+.png}
  \caption{DeepLab v3+ architecture}
  \label{fig:deeplabv3+}
\end{figure}

\section{PosSegNet}
SegNet is a fully convolutional network for semantic segmentation. It was
proposed by Badrinarayanan et al. \cite{badrinarayanan2017segnet} in 2015 and is
structured in two parts: an encoder and a decoder, similar to the one used in
the U-Net. While the encoder is made of convolutional layers followed by max
pooling. The indices of the max pooling are stored to be used in the decoder
during the upsampling. The encoder part is made of the 13 convolutional layers
of the VGG16 network, while in the decoder always 13 layers are used, but they
are made of transposed convolutional layers. Finally, a K-class softmax
classifier is used to predict the class for each pixel.
Also in this case the same positional encoding used in PosPadUNet3D was
introduced, thus named the new network PosSegNet, as it was shown to be
effective in the previous architecture. SegNet was originally proposed for the
segmentation of images, but by changing the 2D operations to 3D ones, it can be
used for the segmentation of 3D volumes. The architecture is shown in figure
\ref{fig:segnet}.
\begin{figure}[h]
  \centering
  \includegraphics[width=0.5\textwidth]{Images/segnet.png}
  \caption{SegNet architecture}
  \label{fig:segnet}
\end{figure}

\section{SwinUNETR3D}
Hatamizadeh et al. \cite{hatamizadeh2022unetr,hatamizadeh2022swin} recently
proposed a novel architectures for 3D medical image segmentation which make use
of the Transformer architecture. The main idea is to replace the convolutions of
U-Net used in the encoding phase with Transformers. Instead of the standard
Vision Transformers, Swin Transformers have been used as they have shown better
performance on images. As Transformers are known to suffer where there is a lack
of data, also a pre-training procedure has been proposed \cite{tang2022self} to
overcome this problem. In the following sections, we will describe SwinUNETR,
starting from the standard Transformer as proposed in the original paper
\emph{"Attention is all you need"} by Vaswani et al.
\cite{vaswani2017attention}, and then we will look at how these ideas have been
moved from Natural Language Processing to Computer Vision.

\subsection{Transformer Architecture}
The transformer architecture was first introduced by Vaswani et al.
\cite{vaswani2017attention} in $2017$ in the field of Natural Language
Processing. The main idea is to replace a recurrent layer with a multi-head
attention mechanism, which is a linear operation that can be computed in
parallel.
The attention mechanism is computed as follows:
\begin{equation}
  \label{eq:attention}
  \text{Attention}(Q,K,V) = \text{softmax}(\frac{QK^T}{\sqrt{d_k}})V
\end{equation}
Where $Q \in \mathbb{R}^{d_k}$ is the query, $K \in \mathbb{R}^{d_k}$ is the key
and $V \in \mathbb{R}^{d_v}$ is the value. If we would like to make it
multi-head, we add more attention layers in parallel and concatenate the
results. The final output of a multi-head attention layer is computed as
follows:
\begin{equation}
  \label{eq:multiheadattention}
  \text{MultiHeadAttention}(Q,K,V) = \text{Concat}(\text{head}_1, \dots,
  \text{head}_h)W^O
\end{equation}
where $\text{head}_i = \text{Attention}(QW_i^Q, KW_i^K, VW_i^V)$.
The projection matrices $W_i^Q \in \mathbb{R}^{d_{\text{model}} \times d_k}$,
$W_i^K \in \mathbb{R}^{d_{\text{model}} \times d_k}$, $W_i^V \in
\mathbb{R}^{d_{\text{model}} \times d_v}$, and $W^O \in \mathbb{R}^{hd_v \times
d_\text{model}}$ are learned parameters. The complete architecture, which also
adds a residual connection, a linear layer, and a normalization layer, is shown
in Figure \ref{fig:transformer}.
\begin{figure}[h]
  \centering
  \includegraphics[width=0.8\textwidth]{transformer.png}
  \caption{Transformer architecture}
  \label{fig:transformer}
\end{figure}

\subsection{Vision Transformer}
While the Transformer architecture has become the de-facto standard for natural
language processing tasks, its applications to computer vision have remained
limited for a few years. One of the reasons is that the complexity of the
attention mechanism is $O(n^2)$, where $n$ is the number of pixels in the image.
This makes it impractical for medium and large-size images. In $2020$,
Dosovitskiy et al. \cite{dosovitskiy2020image} proposed a novel architecture
called Vision Transformer (ViT) which is able to scale up the Transformer
architecture to images. The main idea is to reshape the original image $x \in
\mathbb{R}^{H \times W \times C}$ into a sequence of patches $x_p \in
\mathbb{R}^{P \times P \times C}$, where $H$, $W$, $C$ and $P$ are the height,
width, number of channels and resolution of patches respectively. The
Transformer uses constant latent vector size $D$ through all of its layers, so
we flatten the patches and map to $D$ dimensions with a trainable linear
projection. The outputs of such projections are patch embeddings. An additional
token \texttt{[class]} is added to the sequence, which is used to encode the
class label of the image. The whole architecture is depicted in Figure
\ref{fig:visiontransformer}.

\begin{figure}[ht!]
  \centering
  \includegraphics[width=0.5\textwidth]{Images/ViT.png}
  \caption{Vision Transformer architecture}
  \label{fig:visiontransformer}
\end{figure}

\subsection{Swin Transformer}
Swin Transformer, where the word \emph{Swin} stands for \emph{Shifted Window},
is a novel architecture proposed by Liu et al. \cite{liu2021swin} in $2021$
which has become quite popular in the field of computer vision in these latter
years. The problem that the authors wanted to solve is that the Vision
Transformers are not really well suited for tasks that require dense prediction
at pixel level such as semantic segmentation, so they propose a novel that
constructs hierarchical feature maps and has linear computational complexity to
image size.

The Swin Transformer still relies on patches, but instead of choosing one size
and sticking with it, it first starts with small patches for the first
Transformer layer, then merges them into bigger ones in the deeper Transformer
layers. The first size chosen for the patches is $4 \times 4$ and then it is
projected to a chosen dimensionality $C$, which in the paper has been proposed
to be $96$ for the small model and $192$ for the large model. Then, these
vectors are processed with a Transformer which makes use of a Shifted Window
based Self-Attention, introduced in this paper, instead of the classical
Self-Attention used in ViT. This type of attention has the benefit to be linear
because it is limited to a fixed number of patches $M$, so its complexity will
be $O(M \times N)$ instead of $O(N^2)$ where $N$ is the number of patches.

The output of such Transformer Layer is fed into a Merging Layer which
concatenates the vectors of groups of $2 \times 2$ neighboring patches and fed to
another Linear layer to reduce the dimensionality.
This process is repeated but for each layer, the region where the attention was
limited to is shifted in order to apply the attention to a different set of
patches, which before could not be seen among each other.

As a result, Swin Transformers are suitable for various downstream tasks wherein
the extracted multi-scale features can be leveraged for further processing.
\begin{figure}[ht!]
  \centering
  \includegraphics[width=0.8\textwidth]{Images/SwinTransformer.png}
  \caption{Swin Transformer components}
  \label{fig:swintransformer}
\end{figure}

\subsection{Final Architecture and pre-training}
The final architecture that has been proposed by Hatamizadeh et al. \cite{hatamizadeh2022swin} follows the
structure of U-Net3D, but in the encoder part, SwinTransformers are used instead
of the standard 3D Convolution. Some changes from the original paper have been
made in order to adapt the model to our type of data. From the original volume
of size $168 \times 280 \times 360$, we extract patches of size $96 \times 96
\times 96$. These patches are fed in the SwinUNetr model which is composed of a
4-stages Swin Transformer that acts as the encoder part and a sequence of
transposed convolutions as the decoder. For each stage of the encoder part, a
skip connection is added to the output of the decoder stage with the same shape,
as in the original U-Net architecture. The whole architecture is shown in Figure
\ref{fig:swinunetr}.

\begin{figure}[ht!]
  \centering
  \includegraphics[width=0.8\textwidth]{Images/SwinUNETR.png}
  \caption{SwinUNETR architecture}
  \label{fig:swinunetr}
\end{figure}

As already mentioned, a pre-training phase is performed to overcome the data
hungriness characteristic of Transformers. The pretraining has been performed by
paper authors on a variety of datasets of CT scans of different types have been
used. A total of $5$ public CT datasets, consisting of a total of $5050$
subjects, are used to construct the pre-training dataset. We do not have tried
to perform the same pre-training on CBCT data instead of CT because the time
required was prohibitive.
The datasets used are:
\begin{itemize}
  \item{\textbf{Head \& Neck Squamous Cell Carcinoma (HNSCC)}: A collection
    contains imaging, radiation therapy, and clinical data from 627 head and
    neck squamous cell carcinoma (HNSCC) patients at MD Anderson Cancer Center.
    \cite{grossberg2018imaging}}
  \item{\textbf{Lung Nodule Analysis 2016 (LUNA 16)}: a collection of 888 CT
    scans of lung nodules. \cite{setio2017validation}}
  \item{\textbf{TCIA CT Colonography Trial}: a collection of 825 CT colonography
    scans. \cite{johnson2008accuracy}}
  \item{\textbf{TCIA Covid 19}: a set of 2 datasets for a total of 753 CT scans
    of the chest and annotation about the Covid-19 infection.
    \cite{desai2020chest}}
  \item{\textbf{TCIA LIDC-IDRI}: a collection obtained by the collaboration of
    seven academic centers and eight medical imaging companies. It contains 1018
    cases of diagnostic and lung cancer screening thoracic computed tomography
    (CT) scans with marked-up annotated lesions. \cite{armato2011lung}}
\end{itemize}

\begin{figure}[ht!]
  \centering
  \includegraphics[width=0.4\textwidth]{Images/pretrain-unetr.png}
  \caption{Pre-training phase}
  \label{fig:pretrain-unetr}
\end{figure}
Multiple proxy tasks and formulated with a multi-objective loss function has
been used in that self-supervised phase:
\begin{itemize}
  \item{\textbf{Inpainting}: a transposed convolution layer have been attached
    to the encoder as reconstruction head, a patch is removed from the original
    volume and the model must manage to reconstruct the missing voxels, standard
    L1 loss $$\mathcal{L}_{inpaint} = ||\mathcal{X} - \hat{\mathcal{X}}||_1$$ is
    used, where $\hat{\mathcal{X}}$ is the output of the model and $\mathcal{X}$
    is the ground truth.}
  \item{\textbf{Image rotation}: The rotation prediction task predicts the angle
    categories by which the input sub-volume is rotated. For simplicity, we
    employ $R$ classes of $0$, $90$, $180$, and $270$ rotations along the
    z-axis. An MLP classification head is used for predicting the softmax
    probabilities $\hat{y}_r$ of rotation categories. Given the ground truth
    $y_r$, a cross-entropy loss is used for rotation prediction task:
    $$\mathcal{L}_{rot} = -\sum^R_{r=1}y_r\log(\hat{y}_r)$$}
  \item{\textbf{Contrastive coding}: The rotations used above where also used
    inside a contrastive loss \cite{oord2018representation}. Given a batch of
    augmented sub-volumes, the contrastive coding allows for a better
    representation learning by maximizing the mutual information between
    positive pairs, while minimizing that between negative pairs. A linear layer
    is attached to the encoder to map each augmented sub.volume to a latent
    representation $v$. The cosing similarity is used as the distance
    measurement. The final loss between a pair $v_i$ and $v_j$ is defined as:
    $$\mathcal{L}_{contra} = -\log{\frac{\exp{(sim(v_i,v_j)/t)}}{\sum_k^{2N}1_{k
    \neq i}\exp{(sim(v_i,v_k)/t)}}}$$
    where $t$ is the measurement of normalized temperature scale. $1$ is the
    indicator function evaluating to $1 \iff k \neq i$, and $sim(\cdot, \cdot)$
    denotes the dot product between normalized embeddings.
    The contrastive learning loss function has been shown to strengthens the
    intra-class compactness as well as the inter-class separability.}
\end{itemize}

In conclusion, we minimize the final loss that is the sum of these 3 losses:
$$
\mathcal{L} = \lambda_1\mathcal{L}_{inpaint} + \lambda_2\mathcal{L}_{rot} +
\lambda_3\mathcal{L}_{contra}
$$
Where $\lambda_1$, $\lambda_2$ and $\lambda_3$ are the weights of the 3 losses
and are set to $1$ as experimentally shown to be the optimal value. In the
Figure \ref{fig:pretrain-unetr} this process is visually represented.

\section{Hann window function}
By looking at the whole segmented volume, we noticed that when using
PosPadUNet3D, the segmentation was not very accurate in the borders of each
patch. In order to solve this problem, we decided to extract overlapping patches
from the original volume instead of the non-overlapping method used in the
original paper. When two patches overlap, we would like to compute a weighted
average based on the distance from the center.
TODO: explain what happen in audio.
For this reason, we decided to apply the Hann window function to 3D volumes to
try to reduce edge effects.
\subsection{Hann window function in 1D}
The Hann window function, named after Julius von Hann, has been fist defined in
the field of signal processing \cite{oord2018representation}, thus in 1D, as:
$$
w[n] = 0.5(1 - \cos(\frac{2\pi n}{N})), \quad n = 0, \dots, N
$$
where $x$ is the index of the window and $N$ is the total number of samples.
This function is symmetric, with a maximum value of $1$ in the middle of the
window and a minimum value of $0$ at the edges.
Another interesting property is that the sum of two Hann windows shifted by $\frac{N}{2}$
is equal to a rectangular window of width $N$ and height $1$:
\begin{equation}
  \label{eq:hannisone}
  w[n] + w[n + N] = 0.5(1 - \cos(\frac{2\pi n}{N})) + 0.5(1 - \cos(\frac{2\pi (n + N)}{N} )) = 1
\end{equation}
\subsection{Hann window function in 3D}
To apply the Hann window function in 3D, we simply apply the 1D function to each
dimension of the volume:
\begin{equation}
  w[x, y, z] = w[x]w[y]w[z]
\end{equation}
but we have to handle the case where the patch is close to the border of the
original volume. Here the patch is overlapped only on a part of the volume, thus
the window function must be adapted accordingly. In the 3D case, $27$ different
cases can occur, depending on the position of the patch in the original volume.
In Figure \ref{fig:hann-window} all the $9$ cases in the 2D scenario are shown.
Mathematically, this would require to define many different functions and being
unable to efficiently implement in Python the $N$-dimensional case, with $N$ as
a parameter.

% a 3x3 grid figure
\begin{figure}
  \centering
  \begin{subfigure}{0.5\textwidth}
    \centering
    %%%%%%%%%%%%
    \begin{subfigure}{.32\textwidth}
      \centering
      \includegraphics[width=\textwidth]{Images/hann_window_00.jpg}
    \end{subfigure}
    %%%%%%%%%%%
    \begin{subfigure}{.32\textwidth}
      \centering
      \includegraphics[width=\textwidth]{Images/hann_window_10.jpg}
    \end{subfigure}
    %%%%%%%%%%%
    \begin{subfigure}{.32\textwidth}
      \centering
      \includegraphics[width=\textwidth]{Images/hann_window_20.jpg}
    \end{subfigure}
    %%%%%%%%%%%
    \begin{subfigure}{.32\textwidth}
      \centering
      \includegraphics[width=\textwidth]{Images/hann_window_01.jpg}
    \end{subfigure}
    %%%%%%%%%%%
    \begin{subfigure}{.32\textwidth}
      \centering
      \includegraphics[width=\textwidth]{Images/hann_window_11.jpg}
    \end{subfigure}
    %%%%%%%%%%%
    \begin{subfigure}{.32\textwidth}
      \centering
      \includegraphics[width=\textwidth]{Images/hann_window_21.jpg}
    \end{subfigure}
    %%%%%%%%%%%
    \begin{subfigure}{.32\textwidth}
      \centering
      \includegraphics[width=\textwidth]{Images/hann_window_02.jpg}
    \end{subfigure}
    %%%%%%%%%%%
    \begin{subfigure}{.32\textwidth}
      \centering
      \includegraphics[width=\textwidth]{Images/hann_window_12.jpg}
    \end{subfigure}
    %%%%%%%%%%%
    \begin{subfigure}{.32\textwidth}
      \centering
      \includegraphics[width=\textwidth]{Images/hann_window_22.jpg}
    \end{subfigure}
  \end{subfigure}
  %%%%%%%%%%%
  \caption{Different Hann window functions in 2D based on the position of the extracted patch.}
  \label{fig:hann-window}
\end{figure}

\subsection{Implementation}
The first step is to generate a $N$-dim Hann window, by starting from a $1$-dim
window. Here we use the method provided in the PyTorch library \texttt{torch.hann\_window()}.
To produce a $2$-dim window, is possible to exploit the broadcast property of
tensors by multiplying a $1$-dim window with shape $(N,1)$
with a $1$-dim window with shape $(1,N)$, resulting in a $2$-dim window with
shape $(N,N)$. This process can be repeated to generate up to a $N$-dim window.
The code is shown in Listing \ref{alg:hann-window}. %Listing with title
\begin{lstlisting}[language=Python, caption=Python code to generate a $N$-dim Hann window, label=alg:hann-window]
def _get_hann_window(patch_size):
  hann_window_nd = torch.as_tensor([1])
  # create a n-dim hann window
  for spatial_dim, size in enumerate(patch_size):
    window_shape = np.ones_like(patch_size)
    window_shape[spatial_dim] = size
    hann_window_1d = torch.hann_window(
      size + 2,
      periodic=False,
    )
    hann_window_1d = hann_window_1d[1:-1].view(*window_shape)
    hann_window_nd = hann_window_nd * hann_window_1d
  return hann_window_nd
\end{lstlisting}

Once a $N$-dim Hann window is generated, the next step is to apply it to the
predicted segmnetation and, to avoid small numerical errors, we save the
coefficients used in a tensor of the same shape of the segmentation. These
numerical errors are due to the fact that these Hann windows are discrete and
not continuous, thus the sum of the window coefficients is not always equal to
$1$. Moreover this approach to avoid us to craft a specific window tensor for each
case near the border, which would become impratical for high dimensional cases.
The code is shown in Listing \ref{alg:apply-hann-window}.
\begin{lstlisting}[language=Python, caption=Code to apply Hann window to a batch of patches given their locations, label=alg:apply-hann-window]
hann_window = _get_hann_window(patch_size)
for patch, location in zip(batch, locations):
  i_ini, j_ini, k_ini, i_fin, j_fin, k_fin = location

  patch = patch * _hann_window
  _output_tensor[
      :,
      i_ini:i_fin,
      j_ini:j_fin,
      k_ini:k_fin,
  ] += patch
  _avgmask_tensor[
      :,
      i_ini:i_fin,
      j_ini:j_fin,
      k_ini:k_fin,
  ] += _hann_window
\end{lstlisting}

Once the segmentation for the whole volume is computed, the final step is to
divide the output tensor by the average mask tensor, to obtain the final
segmentation. The code is shown in Listing \ref{alg:final-segmentation}.

\begin{lstlisting}[language=Python, caption=Code to compute the final segmentation, label=alg:final-segmentation]
_output_tensor = _output_tensor / _avgmask_tensor
\end{lstlisting}

This method has been successfully implemented to TorchIO library and is avaible
since version $0.18.83$. The code is available at
\url{https://github.com/fepegar/torchio/blob/main/src/torchio/data/inference/aggregator.py}.

\section{Performance improvements}
The last attempt made aimed to improve the performance of the model without a
excessive drop in the metrics. Being that the face of a person is symmetrical,
we decided to crop each volume in two halves, and flip the right one, then we
trained the model on both halves separately. Moreover, as there is a considerable
space between the two halves, we managed to crop a little bit more of data
while leaving the canals still intact. We also wanted to be able to feed a single volume
into the network without having to extract smaller patches from it. To do so, we
resize the original volume by a factor of $3$. Once a resized volume is fed in
the network, we could resize them back to the original size and compute the
metrics as usual. With this approach, we managed to gain a speedup of $24$ times,
while the DICE score was only slightly affected by a drop of $0.02$.

\section{Results}
All the models have been trained using PyTorch \cite{paszke2019pytorch}, on two
NVIDIA P100 GPUs with 16 GB of memory, over the AImagelab server.
The model proposed by Cipriano et al. achieved the best results among the
different tries and managed to keep the title of the state-of-the-art model for
the task of segmentation of the Inferior Alveolar Canal. However, the
experiments that have been made have obtained comparable metrics thus they might
open a new path to follow to improve such results.

% A table with 

% Chapter 6

\chapter{Results}

\lhead{Chapter 6. \emph{Results}} % This is for the header on each page - perhaps a shortened title
\label{chp:results}
%----------------------------------------------------------------------------------------

\def\:{\hskip0pt} %Definisce un modo veloce per permettere a latex di sillabare correttamente anche parole come 4-connectivity. Il corretto utilizzo è il seguente: 4\:-\:connectivity.

\section{PosPadUNet3D}
This architecture, which is the standard U-Net 3D with few additions, represent the state of the art 


%----------------------------------------------------------------------------------------
%	THESIS CONTENT - APPENDICES
%----------------------------------------------------------------------------------------

\addtocontents{toc}{\vspace{2em}} % Add a gap in the Contents, for aesthetics

\appendix % Cue to tell LaTeX that the following 'chapters' are Appendices

%% Appendix A

\chapter{YACCLAB Results} % Main appendix title

\label{AppendixA} % For referencing this appendix elsewhere, use \ref{AppendixA}

\lhead{Appendix A. \emph{YACCLAB Results}} % This is for the header on each page - perhaps a shortened title

In this Appendix we report the automatic generated charts and tables from \textit{YACCLAB} project. As said before in Chapter~\ref{chp:YACCLAB}, tests were performed on each algorithm included in \textit{YACCLAB}, on all datasets, and on four different environments:

\begin{enumerate}
\item \textit{Windows} \textit{PC} with a i7-4790 \textit{CPU} @ 3.60GHz and \textit{Microsoft Visual Studio} 2013; 
\item Intel Core i7-4980HQ \textit{CPU} @ 2.80GHz running \textit{OS~X} with \textit{XCode} 7.2.1;
\item \textit{Windows} \textit{PC} with a i5-6600 \textit{CPU} @ 3.30GHz and \textit{Microsoft Visual Studio} 2013; 
\item Intel Core 2 Duo-T9600 \textit{CPU} @ 2.80GHz running \textit{OS~X} with \textit{XCode} 7.2.1. 
\end{enumerate}

Every test was repeated 10 times, and for each image the minimum execution time was considered. For a detailed discussion over results or details about the acronyms meaning see Section~\ref{sec:results}.

\clearpage

% TODO aggiornare tabelle LINUX completa
% Averages Lorenzo Windows
\begin{table}
	\centering
	\caption{Average results in ms on a i7-4790 \textit{CPU} @ 3.60GHz with \textit{Windows} and \textit{Microsoft Visual Studio} 2013 (lower is better).}
	\label{tab:appendix_ilb14_averages}
	\begin{tabular}{|l|S[table-format=2.3, table-column-width=1.5cm]|S[table-format=2.3, table-column-width=1.5cm]|S[table-format=2.3, table-column-width=1.5cm]|S[table-format=2.3, table-column-width=1.5cm]|S[table-format=2.3, table-column-width=1.5cm]|}
	\hline
	 & {SAUF} & {DiStefano} & {BBDT} & {CT} & {CCIT}\\
	\hline
	MIRflickr  & 0.611 & 0.658  & 0.261 & 0.746  & 0.295\\
    Hamlet     & 6.493 & 7.451  & 5.057 & 9.849  & 6.051\\
	Tobacco800 & 9.702 & 10.981 & 7.595 & 13.371 & 9.481\\
    Medical    & 3.469 & 3.746  & 1.983 & 4.526  & 2.457\\
    Fingerprints & 0.457 & 0.545 & 0.256 & 0.852 & 0.296\\
	3DPeS      & 0.754 & 0.844  & 0.580 & 1.015  & 0.745\\
	\hline
	\end{tabular}
\vskip 1em   
    \begin{tabular}{|l|S[table-format=2.3, table-column-width=1.5cm]|S[table-format=2.3, table-column-width=1.5cm]|S[table-format=2.3, table-column-width=1.5cm]|S[table-format=2.3, table-column-width=1.5cm]|S[table-format=2.3, table-column-width=1.5cm]|}
	\hline
	& {LSL\_STD} & {CTB} & {SBLA} & {PRED} & {NULL}\\
	\hline
	MIRflickr   & 0.464  & 0.397  & 1.451  & 0.328 & 0.187\\
    Hamlet      & 19.422 & 7.418  & 17.455 & 6.155 & 4.274\\
	Tobacco800  & 29.869 & 12.145 & 27.085 & 9.701 & 6.481\\
    Medical     & 7.564  & 3.289  & 8.506  & 2.592 & 1.710\\
    Fingerprints& 0.462 & 0.359 & 1.146 & 0.314 & 0.203\\
	3DPeS       & 2.139  & 1.046  & 2.226  & 0.802 & 0.538\\
	\hline
	\end{tabular}
\end{table}

%Average Stefano Mac
\begin{table}
	\centering
	\caption{Average results in ms on an Intel Core i7-4980HQ \textit{CPU} @ 2.80GHz running \textit{OS~X} with \textit{XCode} 7.2.1 (lower is better).}
	\label{tab:appendix_stefano_averages}
	\begin{tabular}{|l|S[table-format=2.3, table-column-width=1.5cm]|S[table-format=2.3, table-column-width=1.5cm]|S[table-format=2.3, table-column-width=1.5cm]|S[table-format=2.3, table-column-width=1.5cm]|S[table-format=2.3, table-column-width=1.5cm]|}
	\hline
	 & {SAUF} & {DiStefano} & {BBDT} & {CT} & {CCIT}\\
	\hline
    MIRflickr     & 0.420 & 0.663 & 0.296 & 0.750  & 0.349 \\
	Hamlet        & 5.710 & 6.379 & 3.969 & 8.527  & 4.861 \\
	Tobacco800    & 9.085 & 9.551 & 5.987 & 11.193 & 7.957 \\
	Medical       & 2.573 & 3.339 & 1.704 & 3.897  & 2.053 \\
	Fingerprints  & 0.340 & 0.527 & 0.297 & 0.818  & 0.339 \\
	3DPeS         & 0.519 & 0.640 & 0.416 & 0.767  & 0.495 \\
	
	\hline
	\end{tabular}
\vskip 1em   
    \begin{tabular}{|l|S[table-format=2.3, table-column-width=1.5cm]|S[table-format=2.3, table-column-width=1.5cm]|S[table-format=2.3, table-column-width=1.5cm]|S[table-format=2.3, table-column-width=1.5cm]|S[table-format=2.3, table-column-width=1.5cm]|}
	\hline
	& {LSL\_STD} & {CTB} & {SBLA} & {PRED} & {NULL}\\
	\hline
	MIRflickr     & 0.375 & 0.368 & 1.046 & 0.331 & 0.182   \\
	Hamlet        & 11.798 & 4.612 & 12.613 & 4.297 & 2.996 \\
	Tobacco800    & 23.013 & 7.554 & 19.359 & 6.693 & 4.340 \\
	Medical       & 5.662 & 2.159 & 5.897 & 1.949 & 1.282   \\
	Fingerprints  & 0.301 & 0.307 & 0.866 & 0.289 & 0.181   \\
    3DPeS         & 0.798 & 0.464 & 1.566 & 0.420 & 0.274   \\
	\hline
	\end{tabular}
\end{table}

% Averages Canci Windows
\begin{table}
	\centering
	\caption{Average results in ms on a \textit{Windows} \textit{PC} with a i5-6600 \textit{CPU} @ 3.30GHz and \textit{Microsoft Visual Studio} 2013 (lower is better).}
	\label{tab:appendix_canci_averages}
	\begin{tabular}{|l|S[table-format=2.3, table-column-width=1.5cm]|S[table-format=2.3, table-column-width=1.5cm]|S[table-format=2.3, table-column-width=1.5cm]|S[table-format=2.3, table-column-width=1.5cm]|S[table-format=2.3, table-column-width=1.5cm]|}
	\hline
	 & {SAUF} & {DiStefano} & {BBDT} & {CT} & {CCIT}\\
	\hline
	MIRflickr    & 0.776  & 0.834 & 0.409 & 0.965 & 0.508    \\
    Hamlet       & 7.717  & 8.125 & 5.309 & 9.901 & 6.376    \\
	Tobacco800   & 12.441 & 12.563 & 7.974 & 13.422 & 10.159 \\
    Medical      & 3.955  & 4.116 & 2.244 & 4.632 & 2.826    \\
    Fingerprints & 0.535  & 0.644 & 0.329 & 0.985 & 0.400    \\
	3DPeS        & 1.042  & 0.970 & 0.624 & 1.011 & 0.806    \\
	\hline
	\end{tabular}
\vskip 1em   
    \begin{tabular}{|l|S[table-format=2.3, table-column-width=1.5cm]|S[table-format=2.3, table-column-width=1.5cm]|S[table-format=2.3, table-column-width=1.5cm]|S[table-format=2.3, table-column-width=1.5cm]|S[table-format=2.3, table-column-width=1.5cm]|}
	\hline
	& {LSL\_STD} & {CTB} & {SBLA} & {PRED} & {NULL}\\
	\hline
	MIRflickr    & 1.037 & 0.521 & 1.460 & 0.492 & 0.275\\
    Hamlet       & 17.812 & 5.996 & 15.581 & 6.059 & 3.627\\
	Tobacco800   & 29.661 & 10.029 & 24.937 & 9.828 & 5.656\\
    Medical      & 7.932 & 2.820 & 7.949 & 2.680 & 1.579\\
	Fingerprints & 0.671 & 0.382 & 1.115 & 0.378 & 0.231\\
    3DPeS        & 2.217 & 0.826 & 2.002 & 0.808 & 0.455\\
	\hline
	\end{tabular}
\end{table}

% Averages Federico MAC
\begin{table}
	\centering
	\caption{Average results in ms on an Intel Core 2 Duo-T9600 @ 2.80GHz running \textit{OS~X} with \textit{XCode} 7.2.1 (lower is better).}
	\label{tab:appendix_pro_mid2009_averages}
	\begin{tabular}{|l|S[table-format=2.3, table-column-width=1.5cm]|S[table-format=2.3, table-column-width=1.5cm]|S[table-format=2.3, table-column-width=1.5cm]|S[table-format=2.3, table-column-width=1.5cm]|S[table-format=2.3, table-column-width=1.5cm]|}
	\hline
	 & {SAUF} & {DiStefano} & {BBDT} & {CT} & {CCIT}\\
	\hline
	MIRflickr    & 0.948 & 1.289 & 0.613 & 1.496 & 0.698      \\
    Hamlet       & 14.529 & 13.950 & 10.348 & 20.630 & 15.827 \\
	Tobacco800   & 25.209 & 22.553 & 18.119 & 29.343 & 28.360 \\
    Medical      & 6.041 & 6.745 & 3.931 & 8.769 & 5.700      \\
    Fingerprints & 0.678 & 1.024 & 0.567 & 1.555 & 0.611      \\
	3DPeS        & 1.078 & 1.342 & 0.984 & 1.601 & 1.120      \\
	\hline
	\end{tabular}
\vskip 1em   
    \begin{tabular}{|l|S[table-format=2.3, table-column-width=1.5cm]|S[table-format=2.3, table-column-width=1.5cm]|S[table-format=2.3, table-column-width=1.5cm]|S[table-format=2.3, table-column-width=1.5cm]|S[table-format=2.3, table-column-width=1.5cm]|}
	\hline
	& {LSL\_STD} & {CTB} & {SBLA} & {PRED} & {NULL}\\
	\hline
	MIRflickr    & 1.201 & 0.846 & 2.311 & 0.692 & 0.345\\
	Hamlet       & 35.595 & 16.016 & 24.590 & 13.846 & 6.506\\
	Tobacco800   & 63.384 & 29.564 & 37.672 & 25.560 & 10.758\\
    Medical      & 16.544 & 6.367 & 12.354 & 5.047 & 2.465\\
	Fingerprints & 0.817 & 0.649 & 1.706 & 0.583 & 0.349\\
	3DPeS        & 3.322 & 1.511 & 2.920 & 1.101 & 0.598\\
	\hline
	\end{tabular}
\end{table}

%Averages Lorenzo Windows
\begin{figure*}
\centering
\subfigure[MIRflickr] {
	\includegraphics[width=.48\textwidth]{ilb14_averages_mirflickr}
}
\subfigure[Hamlet] {
	\includegraphics[width=.48\textwidth]{ilb14_averages_hamlet}
}
\subfigure[Tobacco800] {
	\includegraphics[width=.48\textwidth]{ilb14_averages_tobacco800}
}
\subfigure[Medical] {
	\includegraphics[width=.48\textwidth]{ilb14_averages_medical}
}
\subfigure[Fingerprints] {
	\includegraphics[width=.48\textwidth]{ilb14_averages_fingerprints}
}
\subfigure[3DPeS] {
	\includegraphics[width=.48\textwidth]{ilb14_averages_3dpes}
}
\caption{Average run-time tests on a i7-4790 \textit{CPU} @ 3.60GHz with \textit{Windows} and \textit{Microsoft Visual Studio} 2013 (lower is better).}
\label{fig:appendix_ilb14_averages}
\end{figure*}

%Averages Stefano Mac
\begin{figure*}
\centering
\subfigure[MIRflickr] {
	\includegraphics[width=.48\textwidth]{stefano_averages_mirflickr}
}
\subfigure[Hamlet] {
	\includegraphics[width=.48\textwidth]{stefano_averages_hamlet}
}
\subfigure[Tobacco800] {
	\includegraphics[width=.48\textwidth]{stefano_averages_tobacco800}
}
\subfigure[Medical] {
	\includegraphics[width=.48\textwidth]{stefano_averages_medical}
}
\subfigure[Fingerprints] {
	\includegraphics[width=.48\textwidth]{stefano_averages_fingerprints}
}
\subfigure[3DPeS] {
	\includegraphics[width=.48\textwidth]{stefano_averages_3dpes}
}
\caption{Average run-time tests on an Intel Core i7-4980HQ \textit{CPU} @ 2.80GHz running \textit{OS~X} with \textit{XCode} 7.2.1 (lower is better).}
\label{fig:appendix_stefano_averages}
\end{figure*}

%Averages Canci Windows
\begin{figure*}
\centering
\subfigure[MIRflickr] {
	\includegraphics[width=.48\textwidth]{canci_averages_mirflickr}
}
\subfigure[Hamlet] {
	\includegraphics[width=.48\textwidth]{canci_averages_hamlet}
}
\subfigure[Tobacco800] {
	\includegraphics[width=.48\textwidth]{canci_averages_tobacco800}
}
\subfigure[Medical] {
	\includegraphics[width=.48\textwidth]{canci_averages_medical}
}
\subfigure[Fingerprints] {
	\includegraphics[width=.48\textwidth]{canci_averages_fingerprints}
}
\subfigure[3DPeS] {
	\includegraphics[width=.48\textwidth]{canci_averages_3dpes}
}
\caption{Average run-time tests on a i5-6600 \textit{CPU} @ 3.30GHz with \textit{Windows} and \textit{Microsoft Visual Studio} 2013 (lower is better).}
\label{fig:appendix_canci_averages}
\end{figure*}


%Averages Federico Mac
\begin{figure*}
\centering
\subfigure[MIRflickr] {
	\includegraphics[width=.48\textwidth]{pro_mid2009_averages_mirflickr}
}
\subfigure[Hamlet] {
	\includegraphics[width=.48\textwidth]{pro_mid2009_averages_hamlet}
}
\subfigure[Tobacco800] {
	\includegraphics[width=.48\textwidth]{pro_mid2009_averages_tobacco800}
}
\subfigure[Medical] {
	\includegraphics[width=.48\textwidth]{pro_mid2009_averages_medical}
}
\subfigure[Fingerprints] {
	\includegraphics[width=.48\textwidth]{pro_mid2009_averages_fingerprints}
}
\subfigure[3DPeS] {
	\includegraphics[width=.48\textwidth]{pro_mid2009_averages_3dpes}
}
\caption{Average run-time tests on an Intel Core 2 Duo-T9600 @ 2.80GHz with \textit{OS~X} and \textit{XCode} 7.2.1 (lower is better).}
\label{fig:appendix_pro_mid2009_averages}
\end{figure*}

% Density Lorenzo Windows
\begin{figure}
	\centering
	\subfigure[Density test] {
		\includegraphics[width=.59\columnwidth]{ilb14_density_cropped}
        \label{fig:appendix_ilb14_density}
	}
    \subfigure[NULL-Normalized Density test] {
		\includegraphics[width=.59\columnwidth]{ilb14_normalized_density_cropped}
        \label{fig:appendix_ilb14_normalized_density}
	}
	\subfigure[Size test] {
		\includegraphics[width=.59\columnwidth]{ilb14_size_cropped}
        \label{fig:appendix_ilb14_size}
	}
    \caption{Density, Size and NULL-Normalized Density tests on a i7-4790 \textit{CPU} @ 3.60GHz with \textit{Windows} and \textit{Microsoft Visual Studio} 2013 (lower is better).}
    \label{fig:appendix_dsn_ilb14}
\end{figure}

%Density Stefano MAC
\begin{figure}
	\centering
	\subfigure[Density test] {
		\includegraphics[width=.59\columnwidth]{stefano_density_cropped}
        \label{fig:appendix_stefano_density}
	}
    \subfigure[NULL-Normalized Density test] {
		\includegraphics[width=.59\columnwidth]{stefano_size_cropped}
        \label{fig:appendix_stefano_normalized_density}
	}
	\subfigure[Size test] {
		\includegraphics[width=.59\columnwidth]{stefano_size_cropped}
        \label{fig:appendix_stefano_size}
	}
    \caption{Density, Size and NULL-Normalized Density tests on an Intel Core i7-4980HQ \textit{CPU} @ 2.80GHz running \textit{OS~X} with \textit{XCode} 7.2.1 (lower is better).}
    \label{fig:appendix_dsn_stefano}
\end{figure}

%Density Canci Windows
\begin{figure}
	\centering
	\subfigure[Density test] {
		\includegraphics[width=.59\columnwidth]{canci_density_cropped}
        \label{fig:appendix_canci_density}
	}
    \subfigure[NULL-Normalized Density test] {
		\includegraphics[width=.59\columnwidth]{canci_normalized_density_cropped}
        \label{fig:appendix_canci_normalized_density}
	}
	\subfigure[Size test] {
		\includegraphics[width=.59\columnwidth]{canci_size_cropped}
        \label{fig:appendix_canci_size}
	}
    \caption{Density, Size and NULL-Normalized Density tests on a \textit{Windows} \textit{PC} with a i5-6600 \textit{CPU} @ 3.30GHz and \textit{Microsoft Visual Studio} 2013 (lower is better).}
    \label{fig:appendix_dsn_canci}
\end{figure}

%Density Federico Mac
\begin{figure}
	\centering
	\subfigure[Density test] {
		\includegraphics[width=.59\columnwidth]{pro_mid2009_density_cropped}
        \label{fig:appendix_pro_mid2009_density}
	}
    \subfigure[NULL-Normalized Density test] {
		\includegraphics[width=.59\columnwidth]{pro_mid2009_normalized_density_cropped}
        \label{fig:appendix_pro_mid2009_normalized_density}
	}
	\subfigure[Size test] {
		\includegraphics[width=.59\columnwidth]{pro_mid2009_size_cropped}
        \label{fig:appendix_pro_mid2009_size}
	}
    \caption{Density, Size and NULL-Normalized Density tests on an Intel Core 2 Duo-T9600 @ 2.80GHz running \textit{OS~X} with \textit{XCode} 7.2.1 (lower is better).}
    \label{fig:appendix_dsn_pro_mid2009}
\end{figure}

%% Appendix A

\chapter{Publications} % Main appendix title

\label{chp:AppendixB} % For referencing this appendix elsewhere, use \ref{AppendixA}

\lhead{Appendix B. \emph{Publications}} % This is for the header on each page - perhaps a shortened title

The research activity conducted in the context of this thesis has led to some publications in international conferences. These are summarized below and included in the following pages.

\begin{enumerate}
\item C. Grana, \textbf{F. Bolelli}, L. Baraldi, R. Vezzani. ``YACCLAB - Yet Another Connected Components Labeling Benchmark'', \textit{International Conference on Pattern Recognition - ICPR 2016}. (Submitted)
\item C. Grana, L. Baraldi, \textbf{F. Bolelli}. Optimized Connected Components Labeling with Pixel Prediction, \textit{Advanced Concepts for Intelligent Vision Systems - ACIVS 2016}. (Submitted)

\end{enumerate}

\clearpage
\pagestyle{empty}
\includepdf[pages=-,noautoscale=true,offset=85 -100,scale=0.95,pagecommand={}]{Papers/2016-icpr-yacclab.pdf}

\clearpage
\pagestyle{empty}
\includepdf[pages=-,noautoscale=true,offset=85 -100,scale=0.95,pagecommand={}]{Papers/2016-acivs.pdf}

\addtocontents{toc}{\vspace{2em}} % Add a gap in the Contents, for aesthetics

\backmatter

%----------------------------------------------------------------------------------------
%	BIBLIOGRAPHY
%----------------------------------------------------------------------------------------

\label{Bibliography}

\lhead{\emph{Bibliography}} % Change the page header to say "Bibliography"

\bibliographystyle{unsrtnat} % Use the "unsrtnat" BibTeX style for formatting the Bibliography

\bibliography{Bibliography} % The references (bibliography) information are stored in the file named "Bibliography.bib"


\end{document}
